\chapter*{Abstract}
\addcontentsline{toc}{chapter}{Abstract}

This work develops a conditional research programme in which temporal structure is constructed from an explicitly declared relational model before spacetime geometry is introduced. The starting data are relational states and directed transitions carrying a normalized probability distribution, a normalized complex representative on phase-bearing sectors, positive transition weights, an ordering parameter distinct from metric clock time, and an inherited dimensionless normalization. The programme then asks which temporal quantities follow from those inputs by definition, exact identity, dynamical closure, or calibrated physical bridge.

Directed transitions carry both an exact Shannon entropy difference and a geometric phase link. Their combination with a separately typed non-exact path-affinity term is used to construct an auditable temporal transition phase, while closed-cycle identities separate exact state-function changes from genuinely non-exact directional information. Positive transition activity then generates an elapsed measure and a relative lapse before a spacetime metric is introduced. Downstream layers develop NOW support, bifurcation, transport, memory, constructive retrodiction, preregistered physical tests, and finally a temporal-to-spacetime closure interface.

The use of the word \emph{derived} is deliberately restricted: IDT does not claim that physical time follows uniquely from Shannon information alone. Its claims are conditional on the declared relational primitives and domains. Standard results from information theory, geometric phase, nonequilibrium response, quantum-state geometry, general relativity, and horizon thermodynamics are treated as imported background or recovery benchmarks rather than as new predictions.

Formal derivation, computational evidence, and physical claim status are maintained as separate layers. Exact identities, reference-model reconstructions, calibrated physical correspondences, and replicated measured results are not interchangeable evidence classes. The strongest present results are constructive and algebraic within their stated domains; the strongest physical claims remain attached to explicit falsification protocols whose calibration variables and null models are fixed independently of the measured outcome.
